% Figure: temperature-entropy diagram of the worked Rankine steam cycle,
% with the four numbered state points carrying the case-study numbers.
% State 1 saturated liquid at condenser pressure (10 kPa, 45.8 C);
% state 2 compressed liquid after the pump (8 MPa); state 3 superheated
% steam at the boiler exit (8 MPa, 500 C); state 4 after the isentropic
% turbine expansion to 10 kPa, two-phase. Entropy axis carries the
% computed entropies; temperature axis in degrees Celsius. The two
% near-coincident left-limb states (1 and 2) are split apart for legibility.
\begin{figure}[htbp]
\centering
\definecolor{engteal}{RGB}{30,120,120}
\begin{tikzpicture}
\begin{axis}[
  width=12cm, height=8.4cm,
  xlabel={specific entropy $s$ (kJ/(kg$\cdot$K))},
  ylabel={temperature $T$ ($^{\circ}$C)},
  xmin=-0.3, xmax=8.5, ymin=0, ymax=560,
  axis lines=left,
  tick label style={font=\scriptsize},
  label style={font=\small},
  clip=true,
]
% vapour dome for water on T-s (schematic but anchored: s_f and s_g at a
% few temperatures, critical point near 374 C, s = 4.41)
\addplot[engblue, very thick, smooth] coordinates {
  (0.0,0) (0.65,45.8) (2.05,150) (3.03,250) (3.87,330) (4.41,374)
};
\addplot[engblue, very thick, smooth] coordinates {
  (4.41,374) (5.12,330) (5.95,250) (6.59,150) (7.40,90) (8.15,45.8)
};
\addplot[only marks, mark=*, engblack, mark size=1.6pt] coordinates {(4.41,374)};
\node[engblack, font=\scriptsize, anchor=south] at (axis cs:4.41,382) {critical point};
\node[engblue, font=\scriptsize] at (axis cs:3.4,120) {two-phase};
% Rankine cycle state points (case-study numbers):
% 1: sat liquid at 10 kPa, T = 45.8 C, s1 = 0.649
% 2: after pump to 8 MPa, T approx 46 C, s2 approx 0.649 (drawn slightly right)
% 3: superheated 8 MPa, 500 C, s3 = 6.727
% 4: isentropic expansion to 10 kPa, s4 = s3 = 6.727, two-phase, T = 45.8 C
% boiler path: 2 -> up left limb to saturation at 8 MPa (295 C) -> across
% two-phase to sat vapour -> up to superheat at 500 C, state 3.
% turbine: 3 -> 4 vertical (isentropic). condenser: 4 -> 1 horizontal.
\addplot[engred, very thick] coordinates {(0.649,45.8) (0.70,47)}; % pump 1->2
\addplot[engorange, very thick, smooth] coordinates {
  (0.70,47) (1.50,120) (2.40,200) (3.40,295) (4.20,295) (5.30,295) (5.74,295)
};
\addplot[engorange, very thick, smooth] coordinates {
  (5.74,295) (6.10,380) (6.45,450) (6.727,500)
};
% turbine 3 -> 4 isentropic (vertical at s = 6.727)
\addplot[engred, very thick] coordinates {(6.727,500) (6.727,45.8)};
% condenser 4 -> 1 horizontal at 45.8 C
\addplot[engteal, very thick] coordinates {(6.727,45.8) (0.649,45.8)};
% markers
\addplot[only marks, mark=*, engblack, mark size=1.6pt] coordinates {
  (0.649,45.8) (0.70,47) (6.727,500) (6.727,45.8)
};
\node[engblack, font=\scriptsize, anchor=north east] at (axis cs:0.649,45.8) {1,2};
\node[engblack, font=\scriptsize, anchor=south west] at (axis cs:6.727,500) {3};
\node[engblack, font=\scriptsize, anchor=north west] at (axis cs:6.727,45.8) {4};
\node[engorange, font=\scriptsize, anchor=south] at (axis cs:3.6,300) {boiler (8 MPa)};
\node[engred, font=\scriptsize, anchor=west] at (axis cs:6.83,270) {turbine};
\node[engteal, font=\scriptsize, anchor=north] at (axis cs:3.4,42) {condenser (10 kPa)};
\end{axis}
\end{tikzpicture}
\caption[Temperature-entropy diagram of the worked Rankine cycle]{The
worked steam Rankine cycle on the temperature-entropy plane, carrying the
case-study state points. The feed pump raises the saturated condensate
(state~1, $10\,\text{kPa}$, $45.8\,\degC$) to the boiler pressure
(state~2, $8\,\text{MPa}$); the two states sit almost on top of one
another on the left limb because the pump work is small, and the figure
splits them apart for legibility. The boiler adds heat at constant
pressure, warming the liquid to saturation, boiling it across the
two-phase region, and superheating the vapour to state~3
($8\,\text{MPa}$, $500\,\degC$). The turbine expands the steam
isentropically down the vertical line of constant entropy to state~4, a
wet vapour at the condenser pressure with quality $0.804$. The condenser
rejects heat at constant pressure back to state~1. The area enclosed by
the loop is the net work per kilogram; the clockwise sense marks a
work-producing engine.}
\label{fig:vol04ch02:rankine-ts-cycle}
\end{figure}
